% TOP_PROBLEM: {"id": 1231, "title": "The Cap Set Problem", "category_id": 2, "category": {"id": 2, "name": "combinatorics", "display_name": "Combinatorics", "description": "Counting problems, graph theory, discrete structures.", "slug": "combinatorics", "order_index": 2, "created_at": "2026-07-31T15:26:25.671Z"}, "status": "open"}
% FIRST_POSTED: 2026-10-01
% ATTEMPT_STATUS: unresolved
% Imported from: unsolvedmath_additional_500.tex; UnsolvedMath ID 1231
% !TeX program = lualatex
% Compile twice: lualatex 1231.tex
% Self-contained: no external bibliography, images, or \input files.
% Numeric section labels and visible numbers are original UnsolvedMath IDs.
\documentclass[11pt,a4paper]{article}
\usepackage[margin=24mm,headheight=14pt]{geometry}
\usepackage{fontspec}
\setmainfont{Latin Modern Roman}
\setsansfont{Latin Modern Sans}
\setmonofont{Latin Modern Mono}
\usepackage{amsmath,amssymb,mathtools}
\usepackage{unicode-math}
\setmathfont{Latin Modern Math}
\usepackage{microtype}
\usepackage{enumitem,xurl,needspace}
\usepackage[dvipsnames]{xcolor}
\usepackage{hyperref}
\hypersetup{unicode=true,colorlinks=true,linkcolor=MidnightBlue,urlcolor=MidnightBlue,
 pdftitle={UnsolvedMath Problem 1231},pdfauthor={Source-annotated compilation},bookmarksnumbered=true}
\usepackage{bookmark,tocloft,titlesec,fancyhdr}
\setlength{\cftsecnumwidth}{4.2em}
\cftsetpnumwidth{2.5em}
\cftsetrmarg{4em}
\titleformat{\section}{\Large\bfseries\raggedright}{\thesection}{0.7em}{}
\pagestyle{fancy}\fancyhf{}
\fancyhead[L]{\small\sffamily Additional UnsolvedMath statements}
\fancyhead[R]{\small Original catalogue IDs}
\fancyfoot[C]{\thepage}
\setlength{\parindent}{0pt}
\setlength{\parskip}{5pt plus 1pt}
\setlength{\emergencystretch}{3em}
\setcounter{tocdepth}{1}\setcounter{secnumdepth}{1}
\setlist[itemize]{leftmargin=3em,itemsep=4pt,topsep=4pt,parsep=0pt}
\allowdisplaybreaks
\newcommand{\N}{\mathbb{N}}
\newcommand{\Z}{\mathbb{Z}}
\newcommand{\Q}{\mathbb{Q}}
\newcommand{\R}{\mathbb{R}}
\newcommand{\C}{\mathbb{C}}
\newcommand{\F}{\mathbb{F}}
\newcommand{\A}{\mathbb{A}}
\newcommand{\PP}{\mathbb{P}}
\newcommand{\E}{\mathbb{E}}
\newcommand{\eps}{\varepsilon}
\newcommand{\dd}{\,\mathrm d}
\newcommand{\sref}[2]{\hyperlink{source:#1:#2}{[S#2]}}
\newif\ifshowcatalogue
\showcataloguetrue
\newcommand{\cataloguescope}[1]{\ifshowcatalogue
 \paragraph{Statement recorded in the supplied source.}
 #1\fi}
\begin{document}
% UnsolvedMath ID 1231; source catalogue code COMB-010

\setcounter{section}{1230}

\section{The Extremal Cap-Set Size over F\_3}\label{1231}

{\small\sffamily UnsolvedMath ID 1231 · Combinatorics}

\subsection{Definitions and mathematical statement}

\paragraph{Shared notation.}Unless a section says otherwise, $\N=\{1,2,\ldots\}$,
$\Z$ is the ring of integers, $\Q,\R,\C$ are the rational, real, and complex
fields, $[N]=\{1,\ldots,\lfloor N\rfloor\}$, and $\log$ is the natural
logarithm. For real functions of a parameter tending to infinity,
$f\ll g$ and $f=O(g)$ mean $|f|\leq Cg$ eventually for some fixed $C$;
$f\gg g$ reverses this comparison; $f=o(g)$ means $f/g\to0$;
$f\sim g$ means $f/g\to1$; and $f\asymp g$ means both order bounds.
A subscript on $O$ or $\ll$ specifies allowed constant dependence.
The expression $n^{o(1)}$ in an upper bound means growth smaller than
$n^\epsilon$ for each fixed $\epsilon>0$. Local definitions govern any
exception, finite-field convention, or use of the same symbol for another
object. An asymptotic claim is not automatically an effective algorithm.



A cap set is $A\subseteq\F_3^n$ with no three distinct elements $x,y,z$ satisfying $x+y+z=0$. In characteristic three this is equivalent to avoiding a nonconstant three-term arithmetic progression. Set $c(n)=\max\{|A|:A\text{ is a cap set in }\F_3^n\}$. Determine the exact or sharp asymptotic size of $c(n)$, rather than merely an exponential upper bound. \sref{1231}{1} \sref{1231}{2}

\paragraph{Statement recorded in the supplied source.}
What is the maximum size of a cap set in $\mathbb{F}_3^n$? \sref{1231}{1}

\paragraph{Residual task reported by the archive.}

The embedded review (2026-08-17) records: Close the gap between the lower capacity 2.2202 and upper capacity 2.7552, and determine exact maxima beyond the known small dimensions. \sref{1231}{1}

\subsection{Short English statement}

How many vectors over the three-element field can be chosen while avoiding every nontrivial three-term arithmetic progression?

\subsection{Research attempt: capacity, exact small dimensions, and the tensor upper-bound gap}
\textbf{Status.} We establish an exact dimension-two calculation, product amplification, and existence of the exponential growth rate. We do not determine that rate. Ellenberg--Gijswijt [S2], whose abstract was inspected, supplies the external exponential upper bound; the elementary arguments below do not replace its polynomial method.

\subsubsection{1. Products give a well-defined asymptotic target}
If $A\subseteq\mathbb F_3^m$ and $B\subseteq\mathbb F_3^n$ are caps, then $A\times B$ is a cap. Indeed a zero-sum triple projects to a zero-sum triple in each factor. In characteristic three, a zero-sum triple with two equal entries has all three equal. Since the factors admit no three distinct zero-sum entries, each coordinate triple is constant, and so is the original triple. Hence
\[c(m+n)\geq c(m)c(n).\tag{1}\]
Put $c(0)=1$. The numbers $a_n=\log c(n)$ are superadditive and satisfy $0\leq a_n\leq n\log3$. For completeness, fix $k$ and write $n=qk+r$, $0\leq r<k$. Equation (1) gives $a_n/n\geq q a_k/n$, whose lower limit is $a_k/k$. Taking the supremum over $k$ and comparing with the tautological upper bound by that supremum proves
\[\lim_{n\to\infty}c(n)^{1/n}
=\sup_{k\geq1}c(k)^{1/k}=:\gamma.\tag{2}\]
Thus a finite-dimensional improvement can be amplified, but establishing one does not automatically identify the limiting capacity.

\subsubsection{2. The affine plane calculation}
Clearly $c(1)=2$. The four points $\{0,1\}^2$ form a cap, by the product argument. Suppose five points formed a cap in the affine plane $\mathbb F_3^2$. There are four directions, each with three parallel lines partitioning the nine points. In any direction the five chosen points must occupy its lines in sizes $2,2,1$, because no line may contain three. Exactly two unordered pairs of chosen points then have that direction. Every pair determines one direction, so all four directions account for eight pairs. But five points have ten pairs, a contradiction. Hence $c(2)=4$.

The same elementary pairing idea gives a general benchmark. Fix $x\in A$. On the $3^n-1$ vectors other than $x$, the involution $y\mapsto -x-y$ has no fixed points. A cap contains at most one member of each pair, since containing both would give a nontrivial zero-sum triple with $x$. Therefore $|A|\leq(3^n+1)/2$. This is valid but much weaker asymptotically than [S2].

\subsubsection{3. Separating the external tensor theorem from its numerical optimization}
A standard quantitative consequence of the polynomial method in [S2] is
\[c(n)\leq3\sum_{j\leq2n/3}[z^j](1+z+z^2)^n.\tag{3}\]
We use (3) as an external result, not as something proved by (1). If $0<t<1$, positivity of coefficients and $t^j\geq t^{2n/3}$ for $j\leq2n/3$ imply
\[c(n)\leq3\left(\frac{1+t+t^2}{t^{2/3}}\right)^n.\]
The stationary point obeys $t(1+2t)/(1+t+t^2)=2/3$, equivalently $4t^2+t-2=0$. Its positive root $t=(\sqrt{33}-1)/8$ lies in $(0,1)$ and gives an upper capacity less than 2.756. Meanwhile $\{0,1\}^n$ proves $\gamma\geq2$. These are deliberately nonoptimal lower bounds, not a claim to improve the constructions reported in the catalogue.

\subsubsection{4. Verification and residual problem}
The diagnostic exhausts all subsets of the nine-point plane, finding maximum 4, and verifies the product construction in dimensions through six. It also checks the derivative equation and numerical upper constant without using the rounded decimal in a proof. The product proof requires the characteristic-three observation about repeated entries; omitting it would leave a gap in ruling out triples with partially coincident projections. The remaining task is either a construction with matching capacity or an upper bound reaching the true product supremum in (2). Neither finite exhaustive data nor the existence of that limit closes the gap.

\subsection{Sources}

{\small

\begin{itemize}

\item[\hypertarget{source:1231:1}{S1}] ULAM.AI, \emph{UnsolvedMath}, supplied \texttt{problems.json}, record 1231 (COMB-010), statement and background. The embedded literature-review date is 2026-08-17; that is the archive’s review date, not a fresh certification. Dataset landing page: \url{https://huggingface.co/datasets/ulamai/UnsolvedMath}.

\item[\hypertarget{source:1231:2}{S2}] J. S. Ellenberg and D. Gijswijt, On large subsets of F\_q\textasciicircum{}n with no three-term arithmetic progression, Annals of Mathematics 185 (2017). \url{https://arxiv.org/abs/1605.09223}. Bibliographic information imported from the supplied record.

\item[\hypertarget{source:1231:3}{S3}] B. Romera-Paredes et al., Mathematical discoveries from program search with large language models, Nature 625 (2024), 468–475. \url{https://doi.org/10.1038/s41586-023-06924-6}. Bibliographic information imported from the supplied record.

\end{itemize}

}

\label{end:1231}
\end{document}
